\scanpage{194}
\mspage{396}

\uline{$\overline{\mathrm{III}}$}\quad \uline{Calculations in} $\mathrm{Gl}(2,\widehat{\mathbb Z})'$ \uline{and in} $\mathrm{Gl}(2,\widehat{\mathbb Z})$.

For a given $u\in\mathcal M_{0,3}^{!\sim}$, whence the quantities
\[
f,g,h,\alpha_0,\beta_0,\alpha,\beta,\xi,\eta:
\tag{53}
\]
(all in $\widehat{\Pi}_{0,3}$, except $\alpha_0,\alpha$ which are only in $\widetilde{\Pi}_{0,3}^{\wedge}\simeq\widehat{\Pi}_{0,3}.\{1,\tau_\infty\}$) we denote by the \unsure{same letters} \unsure{the|their} images in \deleted{$\mathrm{Gl}(2,\widehat{\mathbb Z})$}
\[
\mathrm{Gl}(2,\widehat{\mathbb Z})'\stackrel{\mathrm{dfn}}{=}\mathrm{Gl}(2,\widehat{\mathbb Z})/\mathord{\pm}1,
\]
they are in fact in
\[
\mathrm{Sl}(2,\widehat{\mathbb Z})'=\mathrm{Sl}(2,\widehat{\mathbb Z})/\mathord{\pm}1,
\]
\unsure{except} $\alpha_0,\alpha$ which are only in $\mathrm{Gl}(2,\widehat{\mathbb Z})^{\unsure{\pm'|\pm1}}$
\[
\det\alpha_0=\det\alpha=\varepsilon_3\in\{\pm1\}
\tag{54}
\]
where $\varepsilon_3$ is defined by
\[
\varepsilon_3\equiv\mu\quad(3),
\tag{55}
\]
\illegible{1} otherwise.
\[
\begin{gathered}
f,g,h,\beta_0,\beta,\xi,\eta\in\mathrm{Sl}(2,\widehat{\mathbb Z})'\quad\text{i.e.}\\
\det f=\det g=\det h=\det\beta_0=\det\beta=\det\xi=\det\eta=1
\end{gathered}
\tag{56}
\]
\margin{N.B. $f$ is determined modulo transformation into $\ell_0^n f$, and thereby $\alpha,\beta$ are subject to the same indeterminacy -- more precisely we have an indeterminacy $(\alpha,\beta,f)\mapsto(\ell_0^n\alpha,\ell_0^n\beta,\ell_0^n f)$, (in addition to that of the signs). In what follows, we \unsure{shall assume} therefore \deleted{\illegible{1}} \inserted{$f,\alpha,\beta$} \uline{chosen} in $\widehat{\Pi}_{0,3}$, and that the images in $\mathrm{Gl}(2,\widehat{\mathbb Z})'$ have been \unsure{lifted} to $\mathrm{Gl}(2,\widehat{\mathbb Z})$\ldots}
\margin{cf. page \unsure{398|348} \unsure{or better|in a better form} formulary below}

(\uline{N.B.} \unsure{We} \deleted{\unsure{function}} $\det(-1)=1$, hence $\det$ factors
\[
\mathrm{Gl}(2,\widehat{\mathbb Z})'\xrightarrow{\det}\widehat{\mathbb Z}^{\times},
\]
to give an exact sequence
\[
1\longrightarrow\mathrm{Sl}(2,\widehat{\mathbb Z})'\unsure{\longrightarrow}\mathrm{Gl}(2,\widehat{\mathbb Z})'\longrightarrow\widehat{\mathbb Z}^{\times}\unsure{\longrightarrow}1
\tag{57}
\]
).

We shall choose representatives
\[
f,\unsure{\alpha|\alpha_0},\unsure{\beta_0|\beta}\in\mathrm{Gl}(2,\widehat{\mathbb Z})
\tag{58}
\]
which is equivalent to the choice of \uline{three} \uline{signs},

\scanpage{195}

and in terms of these choices we define the other quantities by
\[
\deleted{\left\{\begin{array}{ll}
g=\alpha_0\rho(\alpha_0)^{-1},&h=\beta_0\sigma_\infty(\beta_0)^{-1},\quad\text{and}\\
\alpha_0=f^{-1}\unsure{\alpha_0|\alpha},&\beta_0=f^{-1}\unsure{\beta_0|\beta},\\
\xi=\alpha\rho(\alpha)^{-1},&\eta=\beta\rho(\beta)^{-1}
\end{array}\right.}
\]
\[
\left\{\begin{array}{ll}
\xi=\alpha\rho(\alpha)^{-1},&\eta=\beta\deleted{\unsure{\rho}}\inserted{\sigma_\infty}(\beta)^{-1}\\[4pt]
\alpha_0=f^{-1}\alpha,&\beta_0=f^{-1}\beta\\[8pt]
g=\alpha_0\rho(\alpha_0)^{-1}=f^{-1}\xi\rho(f),&h=\beta_0\rho(\beta_0)^{-1}=f^{-1}\eta\rho(f).
\end{array}\right.
\tag{59}
\]

We \unsure{shall note} that $\xi$ and $\eta$, \deleted{\unsure{or}} $g,h$ are independent of the choice of signs that have been made. The essential relation linking $(f,\unsure{\alpha_0|\alpha},\unsure{\beta_0|\beta})$, or rather $\alpha,\beta$ (in addition to the relations (59), which play the role of defining relations) is the ``loop relation'', in matrix form, which takes the form (corresponding to (46) or to (43))
\[
\boxed{\xi=\varepsilon\eta\ell_1^\nu}\qquad\varepsilon\in\{\pm1\}
\tag{60}
\]
where $\varepsilon\deleted{\unsure{\varepsilon_2}}$ is a sign which no longer depends on any choice -- \unsure{we have seen|we have in} \deleted{\illegible{1}} (cf. also formulary below)
\[
\varepsilon=\varepsilon_2\quad
\text{defined in terms }\inserted{\unsure{\text{i.e.}}}\text{ of }\mu\text{ by }\quad
\mu\equiv\varepsilon_2\quad(4),
\tag{61}
\]

\scanpage{196}
\mspage{397}

and where $\ell_1$ is defined in matrix form by
\[
\deleted{\unsure{(62)}}\qquad
\ell_1=\begin{pmatrix}1&0\\2&1\end{pmatrix}
\]
whence
\[
\ell_1^\nu=\begin{pmatrix}1&0\\2\nu&1\end{pmatrix}
=\begin{pmatrix}1&0\\\mu-1&1\end{pmatrix}.
\tag{62}
\]

\uline{Remark}

(Relation with the notations of \S\,44 (page 356 ff, in particular formulary after page 364). In the groups $\widehat{\Pi}_{0,3}$ and $\widetilde{\Pi}_{0,3}^{\wedge}$, the letters $f,g,h$ of loc. cit. \unsure{have|are} the same meaning as here, but not the letters $\alpha,\beta$ etc. I shall distinguish by underlines, in the case of \unsure{deviation} \inserted{between the two notations}, the letters with the \deleted{\unsure{meaning}} meaning \unsure{used|studied} in loc. cit, we find
\[
\uline{\alpha}=\alpha_0,\qquad\uline{\beta}=\beta_0
\tag{i}
\]
from page 357 on we practically always assume $\uline{\beta}=\beta_0=1$, \inserted{(i.e. $u(\sigma_\infty)=\sigma_\infty$)} so that
\[
\boxed{\beta=f\beta_0=f}
\tag{ii}
\]
\deleted{and} and loc. cit (17) gives
\[
\uline{\varphi}=f\sigma_\infty(f)^{-1}=\beta\sigma_\infty(\beta)^{-1}=\eta
\tag{iii}
\]
whereas (\unsure{loc.} cit (14))
\[
\uline{\xi}=\ell_0^\nu\uline{\varphi}=\ell_0^\nu\eta,
\]
comparing with
\[
\xi=\eta\ell_1^\nu\qquad(46)
\]
we find
\[
\deleted{\unsure{(iv)}}\quad\uline{\xi}=\ell_0^\nu\xi\ell_1^{-\nu}
=\ell_0^\nu\xi\rho(\ell_0^\nu)^{-1}.
\tag{iv}
\]
i.e. $\uline{\xi}$ is the $\rho$-transform of $\ell_0^\nu$.
The relation (loc. cit (\unsure{19|29}))
\[
\uline{\xi}=\uline{\alpha}_0\rho(\uline{\alpha}_0)^{-1}
\]
compared with $\xi=\alpha\rho(\alpha)^{-1}$ and with (iv)\deleted{\illegible{1}}, gives

\scanpage{197}

\[
\deleted{(\illegible{1})}\qquad\uline{\alpha}_0=\ell_0^\nu\alpha.
\tag{v}
\]

The formulas \deleted{\illegible{2}} \inserted{(i) to (v)} make it possible to translate the formulary of loc. cit. in terms of the formulary that I am going to make explicit.

We make $\alpha$ and $\beta$ explicit by their matrices
\[
\alpha=\begin{pmatrix}a&b\\c&d\end{pmatrix},\qquad
\beta=\begin{pmatrix}r&s\\t&u\end{pmatrix},
\tag{63}
\]
where $a,b,c,d,r,s,t,u\in\widehat{\mathbb Z}$, with
\[
ad-bc=\deleted{\varepsilon_{\illegible{1}}}\inserted{\delta}=\varepsilon_3,\qquad ru-st=1.
\tag{64}
\]
we shall have
\[
\xi=\alpha\rho(\alpha)^{-1}=\begin{pmatrix}A&B\\C&D\end{pmatrix}
\tag{65}
\]
with (cf. page 342, \unsure{(78)|(18)})
\[
\left\{\begin{array}{ll}
\delta A=(a^2+ab+b^2)-(ac+bd+bc),&\delta B=ac+bd+bc\\[4pt]
\delta C=-(c^2+cd+d^2)+(ac+bd+ad),&\delta D=c^2+cd+d^2
\end{array}\right.,
\tag{66}
\]
\[
\eta=\beta\sigma_\infty(\beta)^{-1}=\begin{pmatrix}R&S\\T&U\end{pmatrix},\quad\text{with}
\tag{67}
\]
\[
R=r^2+s^2,\quad S=T=rt+su,\quad U=t^2+u^2.
\tag{68}
\]
The loop relation (60) takes the form
\[
\begin{aligned}
\begin{pmatrix}A&B\\C&D\end{pmatrix}
&=\varepsilon\begin{pmatrix}R&S\\T&\deleted{\unsure{T}}\inserted{U}\end{pmatrix}
\begin{pmatrix}1&0\\2\nu&1\end{pmatrix}\\[4pt]
&=\varepsilon\begin{pmatrix}R+2\nu S&S\\T+2\nu U&U\end{pmatrix},
\end{aligned}
\tag*{\boxed{(69)}}
\]

\scanpage{198}
\mspage{398}

Let us note the relations (\unsure{consequences} of (66), (67)):
\[
(70)\qquad AD-BC=1\qquad RU-ST=1
\]
\[
(71)\qquad C+D-B=1\qquad T\unsure{=|-}S=1\quad\text{i.e. }T=S,
\]
This \deleted{\unsure{cocycle}}\inserted{loop} relation (69) \deleted{takes} is equivalent to the system of relations
\[
(72)\qquad
\boxed{\begin{gathered}
D=\varepsilon U\\
\mu D=1\\
B=\varepsilon S
\end{gathered}}
\qquad(\text{where }\mu=2\nu+1),
\]
i.e. to the system of equations
\[
(73)\qquad
\left\{\begin{aligned}
c^2+cd+d^2&=\varepsilon(t^2+u^2)=\delta/\mu\\
ac+bd+bc&=\varepsilon(rt+su)\\
&\deleted{\mu(c^2+cd+d^2)=\delta}
\end{aligned}\right.
\]
It is \unsure{from the|from} \deleted{\unsure{the comparison of the first}}\inserted{first two relations} (73) that we derive at once
\[
(74)\qquad\delta\equiv\mu\ (3),\qquad\varepsilon\equiv\mu\ (4)
\]
which characterize the signs $\delta,\varepsilon$ as
\[
(75)\qquad\delta=\varepsilon_3(\mu)=\varepsilon_3,\qquad
\varepsilon=\varepsilon_2(\mu)=\varepsilon_2,
\]
cf. \unsure{Theorem} page 365.

The formula (83) on page 370 takes here the very simple form (valid without conditions other than those \unsure{specified} so far on $\alpha,\beta,f$ --- $f$ moreover not occurring \deleted{\unsure{in}} for the moment in the formulary):

\scanpage{199}

\[
(76)\qquad
\underset{\substack{\Vert\\\varepsilon_3\in\{\pm1\}}}{(\delta)}
\alpha(c\rho+d)=\beta(t\sigma_\infty+u)\tau_\infty^{\alpha_2}
\qquad(\alpha=\alpha_2)
\]
where $\alpha_2=\alpha$ is defined by the condition
\[
(-1)^{\alpha_2}=\varepsilon\quad(=\varepsilon_2).
\]
Making (76) explicit by matrix calculation, we have
\[
\delta\alpha(c\rho+d)
=\delta\begin{pmatrix}a&b\\c&d\end{pmatrix}
\begin{pmatrix}d&c\\-c&c+d\end{pmatrix}
=\deleted{\unsure{\delta}}
\begin{pmatrix}1&B\\0&D\end{pmatrix}
\]
\[
\beta(t\sigma_\infty+u)\tau_\infty^\alpha
=\begin{pmatrix}\unsure{r}&\unsure{s|-s}\\\unsure{t|-t}&u\end{pmatrix}
\begin{pmatrix}u&\unsure{\varepsilon_2t|\varepsilon t}\\-t&\varepsilon u\end{pmatrix}
=\begin{pmatrix}1&\varepsilon S\\0&\varepsilon U\end{pmatrix}
\]
so that (76) reduces to the equations
\[
D=\varepsilon U,\quad B=\varepsilon S,
\]
thus \uline{the loop relation} $\xi=\varepsilon\eta \ell_1^\nu$ (\unsure{60|68}, 69) \uline{is equivalent} to \uline{the conjunction of} (76) and of the relation
\[
(77)\qquad\mu D=1.
\]
This being laid down, we shall set \inserted{(for $(\alpha,\beta)$ fixed so as to satisfy the loop formula)}
\[
\begin{aligned}
(78)\qquad u_{\alpha,\beta}
&\overset{\mathrm{dfn}}{=}(\mu\delta).\alpha(c\rho+d)
=\mu\beta(t\sigma_\infty+u)\tau_\infty^\alpha\\
&=\mu\begin{pmatrix}1&B\\0&D\end{pmatrix}
=\begin{pmatrix}\mu&\mu B\\0&1\end{pmatrix}
\end{aligned}
\]
so that $u_{\alpha,\beta}$ in fact depends only on $\alpha$, or only on $\beta$ and the sign $\alpha_2$ (which, for its part, is not determined by $\beta$, but by $\alpha$ as the value of $\delta(c^2+cd+d^2)\bmod4$\dots); and we shall set, for a triple $(\alpha,\beta,f)$ of elements of $\mathrm{Gl}(2,\widehat{\mathbb Z})$, $\det f=1$, $\unsure{\alpha\text{ and }\beta|\det\beta}$ related by the loop formula

\margin{in fact, $u_{\alpha\beta}$ depends only on $\xi=\alpha\rho(\alpha)^{-1}$, and even only on the image $\unsure{\overline\xi|\xi}$ of $\xi$ in $\unsure{\mathrm{Sl}(2,\widehat{\mathbb Z})|\mathrm{Sl}(2,\widehat{\mathbb Z})'}$, \unsure{the sign being characterized} by the condition $\xi\rho(\xi)\rho^2(\xi)=+1$) i.e. depends only on $\unsure{B,D(=1/\mu\dots)}$. It should be denoted $u_\xi$, cf. page 410.}